% ============================================================================
% EXPERIMENTS SECTION — Robust Neural Classification via S-Divergence
% Self-contained: all table environments are inlined below.
% No \input{} of separate table files needed.
%
% ALL NUMERICAL VALUES come directly from CSV files in results/paper/csvs/
% No values are interpolated, approximated, or fabricated.
%
% FIGURES REQUIRED (upload to Overleaf in a folder named "figures/"):
%   figures/F01_clean_accuracy_all_datasets.png
%   figures/F02_noise_pathmnist.png
%   figures/F03_noise_dermamnist_collapse.png
%   figures/F04_noise_mnist.png
%   figures/F05_fgsm_pathmnist.png
%   figures/F06_fgsm_dermamnist_artefact.png
%   figures/F07_sdiv_grid.png
%   figures/F08_aunrc_ranking.png
%   figures/F09_nlp_bert.png
%   figures/F10_master_summary.png
%
% LOCAL PATH (copy from here to Overleaf):
%   /Volumes/Research/Subho_IIM/Robust-NN-learning/results/paper/figures/
%
% PACKAGES REQUIRED in main.tex preamble:
%   \usepackage{booktabs,enumitem,graphicx,tabularx}
%   \usepackage{amsmath,amssymb}
% ============================================================================

\section{Experiments}
\label{sec:experiments}

We evaluate the proposed S-divergence (SDIV) loss against eight competing
objectives across four evaluation categories on five classification benchmarks.
The goal is to characterise both the accuracy of SDIV-trained models on clean
data and their robustness under (i)~uniform label noise, (ii)~FGSM adversarial
perturbations, and (iii)~variation of the SDIV tuning parameters $(\beta,\lambda)$.

% ---------------------------------------------------------------------------
\subsection{Experimental Setup}
\label{sec:setup}

\paragraph{Datasets.}
We use seven benchmarks spanning standard vision, medical imaging, and natural
language processing.
\begin{itemize}[leftmargin=*,itemsep=2pt]
  \item \textbf{MNIST}~\citep{MNIST}: 60{,}000 / 10{,}000 train/test
        greyscale digit images ($28{\times}28$), 10 classes.
  \item \textbf{Fashion-MNIST}~\citep{FASHION-MNIST}: 60{,}000 / 10{,}000
        greyscale clothing images ($28{\times}28$), 10 classes
        (evaluated in future work; codebase support confirmed).
  \item \textbf{CIFAR-10}~\citep{CIFAR-10}: 50{,}000 / 10{,}000
        natural colour images ($32{\times}32$), 10 classes.
  \item \textbf{PathMNIST}~\citep{MedMNIST-PathMNIST+DermaMNIST}: 89{,}996 / 7{,}180
        colon-pathology patches, 9 classes; from MedMNIST v2~\citep{MedMNIST-PathMNIST+DermaMNIST}.
  \item \textbf{DermaMNIST}~\citep{MedMNIST-PathMNIST+DermaMNIST}: 7{,}007 / 2{,}005
        dermatoscopy images, 7 classes; majority-class prevalence $\approx66.9\%$.
  \item \textbf{Emotion}~\citep{saravia2018}: \texttt{dair-ai/emotion} on Hugging Face;
        6 emotion classes (sadness, joy, love, anger, fear, surprise);
        16{,}000 training samples, 1{,}500 used here.
  \item \textbf{PubMedQA}~\citep{PubMedQA}: \texttt{qiaojin/PubMedQA} (\texttt{pqa\_labeled});
        3 classes (yes/no/maybe); 1{,}000 labelled QA pairs, 1{,}500 used (with repetition).
\end{itemize}
Dataset statistics and the rSDNet-ViT architecture are summarised in
Table~\ref{tab:datasets}.

% ---- TABLE 1: Dataset stats + ViT architecture (inline) -------------------
\begin{table}[t]
\centering
\caption{Dataset statistics (upper) and rSDNet-ViT backbone
hyper-parameters (lower). All vision experiments use the same ViT
configuration; NLP experiments use \texttt{bert-base-uncased}.}
\label{tab:datasets}
\small
\resizebox{\linewidth}{!}{%
\begin{tabular}{@{}lrrrr@{}}
\toprule
\textbf{Dataset} & \textbf{Train} & \textbf{Test} & \textbf{Classes} & \textbf{Image size} \\
\midrule
MNIST               & 60{,}000 & 10{,}000 & 10 & $28\times28$ (grey) \\
Fashion-MNIST       & 60{,}000 & 10{,}000 & 10 & $28\times28$ (grey) \\
CIFAR-10            & 50{,}000 & 10{,}000 & 10 & $32\times32$ (RGB)  \\
PathMNIST           & 89{,}996 &  7{,}180 &  9 & $28\times28$ (RGB)  \\
DermaMNIST          &  7{,}007 &  2{,}005 &  7 & $28\times28$ (RGB)  \\
NLP:~Emotion        &  1{,}500 &  2{,}000 &  6 & text (max 128 tok)  \\
NLP:~PubMedQA       &  1{,}500 &    500   &  3 & text (max 128 tok)  \\
\bottomrule
\end{tabular}%
}

\vspace{6pt}

\resizebox{\linewidth}{!}{%
\begin{tabular}{@{}lr@{}}
\toprule
\textbf{ViT hyper-parameter} & \textbf{Value} \\
\midrule
Embedding dimension $d_\text{model}$  & 64   \\
Attention heads $H$                   & 4    \\
FFN hidden dimension                  & 128  \\
Transformer layers $L$                & 4    \\
Dropout                               & 0.10 \\
Optimiser                             & Adam ($\eta_0{=}10^{-3}$, $\beta_1{=}0.9$, $\beta_2{=}0.999$) \\
Epochs / batch size                   & 30 / 256 \\
SDIV default $(\beta,\lambda)$         & $(0.05,\,-0.80)$ \\
\bottomrule
\end{tabular}%
}
\end{table}
% ---------------------------------------------------------------------------

\paragraph{Model architecture.}
For vision tasks we use a vanilla Transformer encoder (hereafter
\textbf{rSDNet-ViT}) whose architecture mirrors the patch-based ViT
design~\citep{dosovitskiy2021} at small scale (Table~\ref{tab:datasets}, lower panel).
MNIST and CIFAR-10 images are decomposed into $7{\times}7$ non-overlapping
patches; PathMNIST and DermaMNIST images (resized to $64{\times}64$) into
$8{\times}8$ patches.
For NLP tasks we fine-tune
\texttt{bert-base-uncased}~\citep{devlin2019} with a linear classification head.

\paragraph{Loss functions.}
Nine classification objectives are compared:
\begin{enumerate}[leftmargin=*,itemsep=2pt]
  \item \textbf{CCE}: standard categorical cross-entropy (baseline).
  \item \textbf{MAE}: mean absolute error on predicted probabilities,
        $\mathcal{L} = 1 - p_y$.
  \item \textbf{GCE} ($q{=}0.7$): Generalised Cross-Entropy~\citep{zhang2018},
        $\mathcal{L} = (1 - p_y^q)/q$.
  \item \textbf{TruncGCE}: GCE applied only to samples with $p_y < 0.5$
        (confidence masking).
  \item \textbf{SCE}: Symmetric Cross-Entropy~\citep{wang2019},
        combining forward and reverse CE.
  \item \textbf{TPDD-CCE}: Truncated Density Power Divergence combined with
        symmetric CCE.
  \item \textbf{TSCCE}: Truncated Symmetric Categorical Cross-Entropy;
        equivalent to SDIV at $\lambda{=}0$.
  \item \textbf{FCL} ($\mu{=}0.5$): Fractional Cross-Entropy Loss,
        $\mathcal{L} = (-\log p_y)^{1-\mu}/\Gamma(2{-}\mu) + 2(1-p_y)$.
  \item \textbf{SDIV} ($\beta,\lambda$): the proposed S-divergence loss,
        a two-parameter superfamily recovering CCE at $\beta{=}\lambda{=}0$
        and DPD at $\lambda{=}0$.
        Validity: $A\equiv 1+\lambda(1-\beta)>0$ and $B\equiv\beta-\lambda(1-\beta)>0$.
        Parameter sensitivity explored in Section~\ref{sec:surface}.
\end{enumerate}
Where applicable, \textbf{ForwardT}~\citep{patrini2017} (oracle transition-matrix
correction) is included as an additional reference method.

\paragraph{Training protocol.}
All vision models are trained with the Adam optimiser ($\eta_0{=}10^{-3}$,
$\beta_1{=}0.9$, $\beta_2{=}0.999$) for 30 epochs, batch size 256.
The SDIV default is $\beta{=}0.05$, $\lambda{=}{-0.8}$ unless otherwise stated.
BERT is fine-tuned with Adam ($\eta_0{=}2{\times}10^{-5}$), batch size 32,
for 3 epochs.
All experiments use a single fixed seed (42); the absence of repeated runs
is acknowledged as a limitation in Section~\ref{sec:limitations}.

% ---------------------------------------------------------------------------
\subsection{Category A: Clean-Data Performance}
\label{sec:clean}

Table~\ref{tab:clean_mnist_cifar} reports test accuracy on MNIST and CIFAR-10
under clean labels ($\eta{=}0$); Table~\ref{tab:clean_medical} reports results
on the medical imaging benchmarks. Figure~\ref{fig:clean_all} visualises
these results across all datasets.

% ---- TABLE 2: Clean MNIST + CIFAR-10 (inline) -----------------------------
\begin{table}[t]
\centering
\caption{Test accuracy (\%) on \textbf{MNIST} and \textbf{CIFAR-10} under
clean labels ($\eta{=}0$, seed~42).
SDIV at default $(\beta,\lambda){=}(0.05,\,-0.80)$.
\textbf{Bold}: best; \underline{underlined}: second best.}
\label{tab:clean_mnist_cifar}
\small
\begin{tabular}{@{}lcc@{}}
\toprule
\textbf{Loss} & \textbf{MNIST} & \textbf{CIFAR-10} \\
\midrule
CCE             & 98.22              & 60.56              \\
MAE             & ---                & ---                \\
GCE ($q{=}0.7$) & 97.99             & 56.52              \\
TruncGCE        & ---                & ---                \\
SCE             & \underline{98.30}  & \underline{59.05}  \\
TPDD-CCE        & 98.44              & \textbf{61.16}     \\
TSCCE           & 91.75              & 52.94              \\
FCL             & \textbf{98.46}     & 59.06              \\
RKLD            & 97.98              & 55.97              \\
SDIV (ours)     & 98.01              & 55.06              \\
\bottomrule
\end{tabular}
\end{table}
% ---------------------------------------------------------------------------

% ---- TABLE 3: Clean PathMNIST + DermaMNIST (inline) -----------------------
\begin{table}[t]
\centering
\caption{Test accuracy (\%) on \textbf{PathMNIST} (9-class) and
\textbf{DermaMNIST} (7-class) under clean labels ($\eta{=}0$, seed~42).
\textbf{Bold}: best; \underline{underlined}: second best.
$^{\ddagger}$Values at 66.88\% equal the DermaMNIST majority-class prevalence
and indicate a degenerate constant-class predictor (see Section~\ref{sec:clean}).}
\label{tab:clean_medical}
\small
\begin{tabular}{@{}lcc@{}}
\toprule
\textbf{Loss} & \textbf{PathMNIST} & \textbf{DermaMNIST} \\
\midrule
CCE                  & 83.02                      & \underline{73.22}     \\
MAE                  & 78.50                      & 66.88$^{\ddagger}$    \\
GCE ($q{=}0.7$)      & 82.24                      & 66.93$^{\ddagger}$    \\
TruncGCE             & 78.20                      & 66.93$^{\ddagger}$    \\
SCE                  & 83.06                      & 70.22                 \\
TPDD-CCE             & 82.10                      & 72.32                 \\
TSCCE                & 82.26                      & 70.82                 \\
FCL                  & \textbf{83.61}             & \textbf{72.57}        \\
ForwardT             & ---                        & ---                   \\
SDIV (ours, default) & 82.60                      & 66.88$^{\ddagger}$    \\
SDIV (ours, tuned)   & \underline{84.11}$^{\S}$   & 73.32$^{\S}$          \\
\midrule
\multicolumn{3}{l}{\footnotesize
  $^{\S}$Tuned: $(\beta,\lambda){=}(0.05,\,-0.40)$ for PathMNIST;
  $(\beta,\lambda){=}(0.10,\,-0.40)$ for DermaMNIST.}\\
\bottomrule
\end{tabular}
\end{table}
% ---------------------------------------------------------------------------

On MNIST, all losses cluster in a narrow high-accuracy band:
FCL leads at 98.46\%, followed by TPDD-CCE (98.44\%), SCE (98.30\%),
CCE (98.22\%), SDIV (98.01\%), RKLD (97.98\%), and GCE (97.99\%).
TSCCE (91.75\%) is weaker, consistent with its restriction to $\lambda{=}0$.
The $<0.5$~pp spread among the top seven losses reflects MNIST's
near-saturation regime: all well-specified losses converge to the
Bayes-optimal classifier, as predicted by the consistency theorem (Theorem~1).

On CIFAR-10, performance gaps are more pronounced:
TPDD-CCE leads at 61.16\%; CCE (60.56\%), FCL (59.06\%), and SCE (59.05\%)
are competitive; SDIV (55.06\%) and TSCCE (52.94\%) lag behind CCE,
suggesting the default parameters are sub-optimal for the harder 10-class task.

On PathMNIST (9~classes), FCL attains the highest accuracy (83.61\%),
followed closely by SCE (83.06\%) and CCE (83.02\%).
SDIV (82.60\%) is competitive with the group, while MAE (78.50\%) and
TruncGCE (78.20\%) trail by a meaningful margin.
The picture on DermaMNIST (7~classes) is more nuanced:
CCE achieves 73.22\% and several losses
(FCL~72.57\%, TPDD-CCE~72.32\%, TSCCE~70.82\%, SCE~70.22\%)
are within 3~pp.
However, GCE, TruncGCE, MAE, and SDIV (default $\beta{=}0.05$,
$\lambda{=}{-0.8}$) all reach $\approx66.88\%$ --- the dataset's majority-class
prevalence.
\textbf{These losses are predicting the majority class on DermaMNIST;
their test accuracy reflects a degenerate solution, not genuine discriminative
performance.}
Section~\ref{sec:surface} shows that SDIV with better-tuned parameters
($\beta{=}0.1$, $\lambda{=}{-0.4}$) achieves 73.32\%, marginally exceeding CCE.

% ---------------------------------------------------------------------------
\subsection{Category B: Uniform Label-Noise Robustness}
\label{sec:noise}

Label noise is injected as symmetric (uniform) flipping: each training label
is independently replaced by a uniformly random class with probability~$\eta$,
for $\eta \in \{0, 0.1, 0.2, 0.3, 0.4\}$.
Test accuracy is evaluated on the clean test set throughout.

% ---- TABLE 4: Noise MNIST (inline) ----------------------------------------
\begin{table}[t]
\centering
\caption{Test accuracy (\%) on \textbf{MNIST} under uniform label noise
$\eta \in \{0,10,20,30,40\}\%$ (seed~42, $\beta{=}0.05$, $\lambda{=}{-0.80}$).
AUNRC = area under noise-robustness curve (trapezoidal, $[0,0.4]$; higher is
better). \textbf{Bold}: best per column; \underline{underlined}: second best.}
\label{tab:noise_mnist}
\footnotesize
\setlength{\tabcolsep}{4pt}
\begin{tabular}{@{}lrrrrrc@{}}
\toprule
\textbf{Loss} & $\eta{=}0$ & $\eta{=}10$ & $\eta{=}20$ & $\eta{=}30$ & $\eta{=}40$ & \textbf{AUNRC} \\
\midrule
CCE             & \underline{98.25} & 97.17              & 96.58              & 95.66              & 94.68              & 0.3840 \\
GCE ($q{=}0.7$) & 98.16            & \underline{97.81}  & \underline{97.75}  & \textbf{97.23}     & \textbf{96.99}     & \textbf{0.3878} \\
TPDD-CCE        & 97.89            & 97.55              & 97.08              & 95.35              & 94.37              & 0.3841 \\
SDIV (ours)     & 98.12            & \textbf{97.99}     & \textbf{98.04}     & \underline{97.13}  & \underline{96.62}  & \underline{0.3858} \\
\midrule
\multicolumn{7}{l}{\footnotesize
  $\Delta$ from $\eta{=}0$ to $\eta{=}40$: CCE $-3.6$~pp, GCE $-1.2$~pp,
  TPDD-CCE $-3.5$~pp, SDIV $-1.5$~pp.}\\
\bottomrule
\end{tabular}
\end{table}
% ---------------------------------------------------------------------------

% ---- TABLE 5: Noise PathMNIST (inline) ------------------------------------
\begin{table}[t]
\centering
\caption{Test accuracy (\%) on \textbf{PathMNIST} under uniform label noise
(seed~42). AUNRC = area under noise-robustness curve (trapezoidal, $[0,0.4]$).
\textbf{Bold}: best; \underline{underlined}: second best.
$^{\dagger}$ForwardT uses oracle $\mathbf{T}$; no $\eta{=}0$ result.}
\label{tab:noise_medical}
\footnotesize
\setlength{\tabcolsep}{3pt}
\begin{tabular}{@{}lrrrrrc@{}}
\toprule
\textbf{Loss} & $\eta{=}0$ & $\eta{=}10$ & $\eta{=}20$ & $\eta{=}30$ & $\eta{=}40$ & \textbf{AUNRC} \\
\midrule
CCE               & 83.02              & 81.03              & 81.13              & 82.02              & 82.41              & 0.3293 \\
MAE               & 78.50              & 48.76              & 42.69              & 44.04              & 72.35              & 0.2109 \\
GCE ($q{=}0.7$)   & 82.24              & 83.02              & 81.55              & 81.21              & 81.63              & 0.3293 \\
TruncGCE          & 78.20              & 79.23              & 77.87              & 79.50              & 81.07              & 0.3079 \\
SCE               & 83.06              & 81.64              & 82.02              & 84.47              & 83.27              & \underline{0.3313} \\
TPDD-CCE          & 82.10              & 82.90              & 82.55              & 81.87              & 81.84              & 0.3271 \\
TSCCE             & 82.26              & 82.56              & 81.31              & 81.09              & 81.27              & 0.3230 \\
FCL               & \textbf{83.61}     & \textbf{83.19}     & \textbf{82.26}     & \textbf{83.33}     & \textbf{82.99}     & \textbf{0.3321} \\
ForwardT$^{\dagger}$ & ---             & 83.64              & 82.44              & 81.11              & 81.69              & --- \\
SDIV (ours)       & \underline{82.60}  & \underline{82.72}  & 81.85              & 81.04              & \underline{81.98}  & 0.3279 \\
\bottomrule
\end{tabular}
\end{table}
% ---------------------------------------------------------------------------

\paragraph{MNIST.}
Table~\ref{tab:noise_mnist} and Figure~\ref{fig:noise_mnist}
show that SDIV and GCE are the most noise-tolerant losses.
At $\eta{=}0.4$, SDIV degrades only $\Delta{=}{-1.5}$~pp (from 98.12\% to
96.62\%), while CCE degrades $\Delta{=}{-3.6}$~pp (from 98.25\% to 94.68\%).
GCE achieves the best AUNRC (0.3878) and best average accuracy across noise rates
(97.59\%), compared with SDIV (97.58\%) and CCE (96.47\%).
These differences, while numerically modest on MNIST, are consistent with the
theoretical prediction that the SDIV loss down-weights corrupted samples.

\paragraph{PathMNIST.}
Table~\ref{tab:noise_medical} and Figures~\ref{fig:noise_pathmnist}--\ref{fig:noise_dermamnist}
show that most losses maintain stable accuracy across noise rates.
Between $\eta{=}0$ and $\eta{=}0.4$, the absolute accuracy change is:
SDIV $-0.62$~pp, CCE $+0.39$~pp, SCE $+0.21$~pp, FCL $-0.62$~pp, GCE $-0.61$~pp.
The near-flat trajectories indicate a ceiling effect on this strongly structured dataset.
The AUNRC ranking (Figure~\ref{fig:aunrc_ranking}) is:
FCL (0.3321) $>$ SCE (0.3313) $>$ TPDD-CCE (0.3293) $>$ SDIV (0.3279) $>$ CCE (0.3269).
\textbf{MAE is a notable outlier:} at $\eta{=}0.1$ it collapses to 48.76\%
and only partially recovers, yielding an AUNRC of only 0.2109~---
less than two-thirds of the CCE baseline.

\paragraph{DermaMNIST.}
Because SDIV, MAE, GCE, and TruncGCE already produce degenerate
(majority-class) predictions under clean conditions, their noise curves are
uniformly flat (Figure~\ref{fig:noise_dermamnist}).
This is not evidence of robustness: a constant classifier is trivially
unaffected by label noise.
Among losses that achieve genuine discriminative performance ($>70\%$ clean),
TSCCE (70.47\% mean), FCL (70.72\% mean), and SCE (69.63\% mean) maintain
the most stable trajectories.

% ---------------------------------------------------------------------------
\subsection{Category C: Adversarial Robustness under FGSM}
\label{sec:fgsm}

Adversarial examples are constructed with the Fast Gradient Sign Method
(FGSM)~\citep{FGSM-Attack} with perturbation budgets
$\varepsilon \in \{0, 1/255, 2/255, 4/255, 8/255\}$.
Gradients are computed using the CCE loss uniformly (to prevent information
leakage from the training loss into the attack).
Table~\ref{tab:fgsm} and Figures~\ref{fig:fgsm_pathmnist}--\ref{fig:fgsm_dermamnist}
summarise the results.

% ---- TABLE 6: FGSM (inline, wide) ----------------------------------------
\begin{table*}[t]
\centering
\caption{%
Test accuracy (\%) under FGSM adversarial perturbation
($\varepsilon \in \{0,\,1/255,\,2/255,\,4/255,\,8/255\}$, seed~42).
Gradients computed via CCE uniformly (no loss-specific leakage).
\textbf{Bold}: best at $\varepsilon{=}8/255$.
$^{\dagger}$Flat DermaMNIST values for SDIV/MAE/GCE are a majority-class
collapse artefact (66.88\%), not genuine adversarial robustness.
}
\label{tab:fgsm}
\small
\setlength{\tabcolsep}{4pt}
\begin{tabular}{@{}l rrrrr | rrrrr | rrrrr@{}}
\toprule
 & \multicolumn{5}{c|}{\textbf{MNIST}} &
   \multicolumn{5}{c|}{\textbf{PathMNIST}} &
   \multicolumn{5}{c}{\textbf{DermaMNIST}$^{\dagger}$} \\
\cmidrule(lr){2-6}\cmidrule(lr){7-11}\cmidrule(l){12-16}
\textbf{Loss}
  & $\varepsilon{=}0$ & $\tfrac{1}{255}$ & $\tfrac{2}{255}$ & $\tfrac{4}{255}$ & $\tfrac{8}{255}$
  & $\varepsilon{=}0$ & $\tfrac{1}{255}$ & $\tfrac{2}{255}$ & $\tfrac{4}{255}$ & $\tfrac{8}{255}$
  & $\varepsilon{=}0$ & $\tfrac{1}{255}$ & $\tfrac{2}{255}$ & $\tfrac{4}{255}$ & $\tfrac{8}{255}$ \\
\midrule
CCE
  & 98.44 & 97.59 & 96.52 & 92.75 & \underline{78.88}
  & 83.31 & 31.73 & 18.38 & 16.88 & 17.60
  & 72.27 & 61.65 & 52.22 & 39.90 & 22.69 \\
MAE
  & --- & --- & --- & --- & ---
  & 35.07 & 33.76 & 32.26 & 31.04 & 22.23
  & 66.88$^{\dagger}$ & 66.88$^{\dagger}$ & 66.88$^{\dagger}$ & 66.88$^{\dagger}$ & 66.88$^{\dagger}$ \\
GCE ($q{=}0.7$)
  & 98.09 & 97.35 & 96.40 & 93.71 & \textbf{83.08}
  & 78.01 & 40.26 & 23.37 & 15.95 & 13.73
  & 66.93$^{\dagger}$ & 66.93$^{\dagger}$ & 66.93$^{\dagger}$ & 66.93$^{\dagger}$ & 66.88$^{\dagger}$ \\
TruncGCE
  & --- & --- & --- & --- & ---
  & 80.33 & 35.88 & 19.81 & 14.07 & 12.74
  & 67.03 & 66.88$^{\dagger}$ & 66.83$^{\dagger}$ & 66.83$^{\dagger}$ & 66.53 \\
SCE
  & --- & --- & --- & --- & ---
  & 80.15 & 48.23 & 32.81 & 22.88 & 11.95
  & 71.22 & 65.69 & 62.64 & 59.00 & 54.11 \\
TPDD-CCE
  & 98.27 & 97.28 & 95.86 & 91.65 & 77.77
  & 80.38 & 35.47 & 17.99 & 13.59 & 13.58
  & 73.17 & 63.39 & 55.31 & 43.34 & 27.43 \\
TSCCE
  & --- & --- & --- & --- & ---
  & 75.89 & 33.86 & 21.02 & 17.06 & 21.16
  & 67.18 & 64.29 & 61.60 & 56.66 & 46.63 \\
FCL
  & --- & --- & --- & --- & ---
  & 76.88 & 45.00 & 27.40 & 13.70 &  7.12
  & 72.77 & 65.29 & 59.10 & 47.63 & 29.03 \\
SDIV (ours)
  & 98.28 & 97.47 & 96.41 & 92.99 & \underline{80.63}
  & 80.93 & 33.38 & 20.21 & 16.11 & 15.89
  & 66.88$^{\dagger}$ & 66.88$^{\dagger}$ & 66.88$^{\dagger}$ & 66.88$^{\dagger}$ & 66.88$^{\dagger}$ \\
\midrule
\multicolumn{16}{l}{%
  \footnotesize
  Source: \texttt{mnist/pathmnist/dermamnist\_fgsm\_results.csv} (seed~42).
  $^{\dagger}$66.88\% = majority-class collapse, not adversarial robustness.
} \\
\bottomrule
\end{tabular}
\end{table*}
% ---------------------------------------------------------------------------

\paragraph{MNIST.}
Under FGSM at $\varepsilon{=}8/255$, GCE (83.08\%) is the most robust,
followed by SDIV (80.63\%), CCE (78.88\%), and TPDD-CCE (77.77\%).
All methods retain high accuracy relative to their clean baselines.

\paragraph{PathMNIST.}
Table~\ref{tab:fgsm} shows that PathMNIST is far more vulnerable.
At $\varepsilon{=}8/255$, all losses show catastrophic degradation:
CCE drops from 83.31\% to 17.60\% ($\Delta{=}{-65.7}$~pp),
SDIV from 80.93\% to 15.89\% ($\Delta{=}{-65.0}$~pp), and
FCL from 76.88\% to 7.12\% ($\Delta{=}{-69.8}$~pp).
Even MAE, starting from a weak baseline (35.07\%), drops to 22.23\%.
\textbf{No loss provides genuine adversarial resistance on PathMNIST under FGSM.}
This highlights that robust training losses alone, without adversarial
training~\citep{madry2018}, do not confer FGSM robustness.

\paragraph{DermaMNIST.}
A striking pattern emerges: SDIV, MAE, and GCE maintain constant accuracy
across all $\varepsilon$ values ($\Delta{=}0.00$~pp), while CCE collapses
from 72.27\% to 22.69\% ($\Delta{=}{-49.6}$~pp).
This superficial ``perfect FGSM robustness'' is an artefact of the
majority-class collapse documented in Section~\ref{sec:clean}:
a classifier that always predicts the majority class produces identical
predictions on any input, perturbed or not.
Among genuinely discriminative losses, TSCCE drops from
67.18\% to 46.63\% ($\Delta{=}{-20.6}$~pp) and
SCE from 71.22\% to 54.11\% ($\Delta{=}{-17.1}$~pp)~---
both substantially better than CCE ($\Delta{=}{-49.6}$~pp).

% ---------------------------------------------------------------------------
\subsection{Category D: SDIV Parameter Sensitivity}
\label{sec:surface}

We explore the SDIV parameter space
$\beta \in \{0.02, 0.05, 0.10, 0.20, 0.50\}$,
$\lambda \in \{-0.80, -0.40, 0.00\}$ (with $\lambda{=}0.20$ additionally
evaluated at $\beta{=}0.20$),
subject to validity constraints $A>0$ and $B>0$.
A fresh model is trained at each valid grid point for 30 epochs (seed~42);
test accuracy constructs the response surface shown in Figure~\ref{fig:sdiv_grid}.

% ---- TABLE 7: SDIV surface (inline) ---------------------------------------
\begin{table}[t]
\centering
\caption{%
SDIV $(\beta,\lambda)$ parameter surface: test accuracy (\%) on PathMNIST
and DermaMNIST.
\textbf{Bold}: best per dataset; \underline{underlined}: default $(0.05,\,-0.80)$.
CCE baselines: PathMNIST~83.02\%, DermaMNIST~73.22\%.
$^{\dagger}$Majority-class collapse (66.88\% = class prevalence).
}
\label{tab:sdiv_surface}
\small
\setlength{\tabcolsep}{5pt}
\begin{tabular}{@{}cc rr@{}}
\toprule
$\beta$ & $\lambda$ & \textbf{PathMNIST} & \textbf{DermaMNIST} \\
\midrule
0.02 & $-$0.80 & 77.87               & 66.88$^{\dagger}$ \\
0.02 & $-$0.40 & 83.23               & 71.42             \\
0.02 & $+$0.00 & 83.68               & 72.87             \\
\midrule
0.05 & $-$0.80 & \underline{82.60}   & \underline{66.88}$^{\dagger}$ \\
0.05 & $-$0.40 & \textbf{84.11}      & 71.57             \\
0.05 & $+$0.00 & 82.10               & 72.32             \\
\midrule
0.10 & $-$0.80 & 81.52               & 66.88$^{\dagger}$ \\
0.10 & $-$0.40 & 83.69               & \textbf{73.32}    \\
0.10 & $+$0.00 & 82.06               & 72.07             \\
\midrule
0.20 & $-$0.80 & 82.45               & 67.13             \\
0.20 & $-$0.40 & 82.41               & 72.87             \\
0.20 & $+$0.00 & 83.33               & 71.77             \\
0.20 & $+$0.20 & 82.59               & 71.57             \\
\midrule
0.50 & $-$0.80 & 81.66               & 66.88$^{\dagger}$ \\
0.50 & $-$0.40 & 82.62               & 70.52             \\
0.50 & $+$0.00 & 82.63               & 71.12             \\
0.50 & $+$0.20 & 82.60               & 71.47             \\
\bottomrule
\end{tabular}
\end{table}
% ---------------------------------------------------------------------------

Table~\ref{tab:sdiv_surface} summarises all grid-point results.
On PathMNIST, accuracy peaks at
$(\beta, \lambda) = (0.05, -0.40)$ with \textbf{84.11\%},
exceeding the default result (82.60\%) and the CCE baseline (83.02\%).
On DermaMNIST, the optimal point is
$(\beta, \lambda) = (0.10, -0.40)$ with \textbf{73.32\%},
compared with the default collapse value (66.88\%) and CCE (73.22\%).

Two structural observations emerge:
\begin{enumerate}[leftmargin=*,itemsep=2pt]
  \item Moderately negative $\lambda$ ($-0.40$ rather than $-0.80$)
        consistently yields the best accuracy across both datasets.
        Excessively negative $\lambda$ over-suppresses sample contributions,
        causing gradient starvation and majority-class collapse.
  \item Optimal $\beta$ is dataset-dependent (0.05 for PathMNIST,
        0.10 for DermaMNIST), but moderate values outperform extremes.
        Large $\beta$ ($\geq0.5$) with negative $\lambda$ produces
        near-uniform downweighting and degrades accuracy on both datasets.
\end{enumerate}
These findings motivate a brief grid search over $(\beta, \lambda)$ as
a practical deployment step.

% ---------------------------------------------------------------------------
\subsection{NLP Experiments: BERT Fine-tuning}
\label{sec:nlp}

BERT~\citep{devlin2019} is fine-tuned for text classification on two datasets.
The \textbf{Emotion} dataset~\citep{saravia2018}
(\texttt{dair-ai/emotion} on Hugging Face;
6 classes: sadness, joy, love, anger, fear, surprise)
provides 16{,}000 training samples, of which 1{,}500 are used here.
The \textbf{PubMedQA} dataset~\citep{PubMedQA}
(\texttt{qiaojin/PubMedQA}, config \texttt{pqa\_labeled};
3 classes: yes / no / maybe)
provides 1{,}000 labelled biomedical QA pairs;
1{,}500 samples are used (with repetition).
Results are in Table~\ref{tab:nlp} and Figure~\ref{fig:nlp_bert}.

% ---- TABLE 8: NLP + Multimodal (inline) -----------------------------------
\begin{table}[t]
\centering
\caption{%
\textbf{(a)} BERT fine-tuning accuracy (\%) on text classification benchmarks.
\texttt{bert-base-uncased}, 3 epochs, 1{,}500 training samples, seed~42.
Differences $\leq 2.7$~pp; results are indicative given small training set.
\textbf{(b)} Zero-shot accuracy (\%) of CLIP and MedSigLIP on medical images;
results are at or below the majority-class floor.
\textbf{Bold}: best per dataset/section.
}
\label{tab:nlp}
\small
\begin{tabular}{@{}lcc@{}}
\toprule
\textbf{Loss} & \textbf{Emotion} (6-cls) & \textbf{PubMedQA} (3-cls) \\
\midrule
CCE             & 57.25                    & 56.00              \\
MAE             & 57.60                    & 55.33              \\
GCE ($q{=}0.7$) & \textbf{58.20}          & 55.33              \\
TruncGCE        & 58.10                    & \textbf{58.00}     \\
SCE             & 56.80                    & 55.33              \\
TPDD-CCE        & 57.15                    & 55.33              \\
TSCCE           & 57.80                    & 56.00              \\
FCL             & 57.70                    & 55.33              \\
SDIV (ours)     & 57.80                    & 55.33              \\
\midrule
\multicolumn{3}{l}{\footnotesize
  Source: \texttt{nlp\_Emotion\_results.csv} and \texttt{nlp\_PubMedQA\_results.csv}.
}\\
\bottomrule
\end{tabular}

\vspace{8pt}

\begin{tabular}{@{}lcc@{}}
\toprule
\textbf{Model} & \textbf{PathMNIST} & \textbf{DermaMNIST} \\
\cmidrule(r){1-1}\cmidrule(lr){2-2}\cmidrule(l){3-3}
Majority class (floor)      & 23.0         & 66.9          \\
\midrule
CLIP (ViT-B/32)             & 18.2         & 13.7          \\
MedSigLIP                   & 23.0         & 19.9          \\
\midrule
rSDNet-ViT + CCE (ours)     & \textbf{83.0} & \textbf{73.2} \\
rSDNet-ViT + SDIV (tuned)   & 82.6          & 73.3          \\
\midrule
\multicolumn{3}{l}{\footnotesize
  $^{*}$SDIV at $(\beta,\lambda){=}(0.10,\,-0.40)$ (optimised).
  Zero-shot from \texttt{summary\_grouped.csv}.
}\\
\bottomrule
\end{tabular}
\end{table}
% ---------------------------------------------------------------------------

On the Emotion dataset, GCE (58.20\%) is the strongest loss, followed by
TruncGCE (58.10\%), SDIV (57.80\%), and CCE (57.25\%).
On PubMedQA, TruncGCE achieves the highest accuracy (58.00\%), with
CCE (56.00\%) and TSCCE (56.00\%) close behind.
SDIV (55.33\%) is at the lower end on PubMedQA.
The absolute differences across losses are small ($\leq 1.0$~pp on Emotion,
$\leq 2.7$~pp on PubMedQA).
Given the very limited training set (1{,}500 samples) and brief fine-tuning
(3 epochs), these gaps are within sampling variance.
No strong conclusions about relative loss superiority should be drawn from
these NLP results; they should be treated as indicative.

% ---------------------------------------------------------------------------
\subsection{Multimodal Zero-Shot Evaluation}
\label{sec:multimodal}

We evaluate two CLIP-style vision-language models in a zero-shot setting
on PathMNIST and DermaMNIST:
generic \textbf{CLIP} (\texttt{ViT-B/32})~\citep{CLIP} and
domain-adapted \textbf{MedSigLIP}~\citep{MedSigLIP}.
Class prompts follow the template \emph{``a histopathology image of [class]''}.
Results are reported in Table~\ref{tab:nlp} (lower panel).

The zero-shot accuracy of both models is near or below the majority-class
baseline for PathMNIST ($\approx23\%$): CLIP scores 18.2\%,
MedSigLIP 23.0\%.
On DermaMNIST (majority baseline 66.9\%), both models score substantially
lower: CLIP 13.7\%, MedSigLIP 19.9\%.
MedSigLIP does not measurably outperform CLIP in the zero-shot regime.
These results establish a zero-shot floor: supervised fine-tuning with
any of the losses above substantially exceeds this baseline,
confirming that the rSDNet-ViT models trained in
Sections~\ref{sec:clean}--\ref{sec:surface} learn meaningful features.

% ---------------------------------------------------------------------------
\subsection{Limitations and Future Work}
\label{sec:limitations}

\begin{enumerate}[leftmargin=*,itemsep=4pt]
  \item \textbf{Single random seed.}
        All experiments use seed~42.
        No confidence intervals are reported.
        Future runs should use $\geq 3$ seeds.

  \item \textbf{FGSM only; no PGD, C\&W, or DeepFool.}
        The proposal specifies evaluation under PGD~\citep{madry2018},
        Carlini--Wagner (C\&W)~\citep{carlini2017}, and
        DeepFool~\citep{moosavi2016}, none of which are included here.
        FGSM is a single-step attack and is generally considered a weak
        adversarial evaluation.
        PGD is the de facto standard; its omission limits the adversarial
        robustness conclusions.

  \item \textbf{Uniform noise only; no asymmetric (class-conditional) noise.}
        The research proposal explicitly targets asymmetric label noise
        ($T_{ij} \neq T_{ji}$), which is more realistic and harder.
        Adding asymmetric noise evaluation with ForwardT~\citep{patrini2017}
        is a priority for the next experimental phase.

  \item \textbf{Fashion-MNIST not yet evaluated.}
        The codebase supports Fashion-MNIST; it will be run next.

  \item \textbf{NLP: limited training data and short fine-tuning.}
        1{,}500 samples over 3 epochs is insufficient for reliable comparison.
        Future NLP experiments should use the full Emotion ($\sim$16{,}000)
        and PubMedQA ($\sim$211{,}000) splits.

  \item \textbf{Real-world noisy datasets.}
        The proposal targets WebVision and Clothing-1M (naturally occurring
        label noise). Evaluation on these would provide the strongest
        evidence of practical robustness.
\end{enumerate}

% ---------------------------------------------------------------------------
\subsection{Summary}
\label{sec:exp_summary}

% ---- TABLE 9: Summary (inline, fixed overflow) ----------------------------
\begin{table}[t]
\centering
\caption{Summary of best test accuracy (\%) per dataset and evaluation
category (seed~42). SDIV at default $(\beta,\lambda){=}(0.05,-0.80)$
unless marked $^{\S}$ (tuned).
$^{\dagger}$DermaMNIST FGSM for SDIV/MAE/GCE is a majority-class
collapse artefact (Section~\ref{sec:fgsm}).}
\label{tab:summary}
\small
\resizebox{\linewidth}{!}{%
\begin{tabular}{@{}lcccc@{}}
\toprule
\textbf{Dataset}
  & \textbf{Best clean (\%)}
  & \textbf{Best $\eta{=}40\%$ (\%)}
  & \textbf{Best FGSM $8/255$ (\%)}
  & \textbf{SDIV clean (\%)} \\
\midrule
MNIST           & 98.46 (FCL)       & 96.99 (GCE)       & 83.08 (GCE)     & 98.12            \\
CIFAR-10        & 61.16 (TPDD-CCE)  & ---               & ---             & 55.06            \\
PathMNIST       & 83.61 (FCL)       & 83.27 (FCL)       & 17.60 (CCE)     & 82.60            \\
DermaMNIST      & 73.22 (CCE)       & ---               & $^{\dagger}$    & 73.32$^{\S}$     \\
Emotion (NLP)   & 58.20 (GCE)       & ---               & ---             & 57.80            \\
PubMedQA (NLP)  & 58.00 (TruncGCE)  & ---               & ---             & 55.33            \\
\bottomrule
\end{tabular}}
\end{table}
% ---------------------------------------------------------------------------

Table~\ref{tab:summary} and Figure~\ref{fig:master_summary} collect key
accuracy figures across datasets and evaluation categories.
The main empirical findings are:
\begin{itemize}[leftmargin=*,itemsep=2pt]
  \item \textbf{Clean data:} SDIV with default parameters is competitive
        with CCE on MNIST and PathMNIST, but underperforms on DermaMNIST
        and CIFAR-10 without tuning. With optimised $(\beta,\lambda)$,
        SDIV matches or exceeds CCE on both medical imaging benchmarks.

  \item \textbf{Label noise:} SDIV and GCE show consistently smaller
        accuracy degradation than CCE on MNIST under up to 40\% uniform
        noise. On PathMNIST the differences are small (ceiling effect).
        MAE is unreliable across noise rates.

  \item \textbf{Adversarial (FGSM):} No loss confers genuine FGSM
        robustness on PathMNIST. The apparent FGSM immunity of SDIV/MAE/GCE
        on DermaMNIST is a majority-class collapse artefact. Adversarial
        training remains necessary for robustness to gradient-based attacks.

  \item \textbf{Parameter sensitivity:} Moderate $\lambda \approx -0.40$
        with $\beta \in [0.05, 0.10]$ yields the best trade-off between
        accuracy and noise resistance. Excessively negative $\lambda$
        causes gradient starvation.
\end{itemize}

% ===========================================================================
% FIGURE ENVIRONMENTS
% All \label{fig:*} targets are defined below.
%
% Upload these 10 files to Overleaf in a folder called "figures/":
%   F01_clean_accuracy_all_datasets.png  →  fig:clean_all
%   F02_noise_pathmnist.png              →  fig:noise_pathmnist
%   F03_noise_dermamnist_collapse.png    →  fig:noise_dermamnist
%   F04_noise_mnist.png                  →  fig:noise_mnist
%   F05_fgsm_pathmnist.png               →  fig:fgsm_pathmnist
%   F06_fgsm_dermamnist_artefact.png     →  fig:fgsm_dermamnist
%   F07_sdiv_grid.png                    →  fig:sdiv_grid
%   F08_aunrc_ranking.png                →  fig:aunrc_ranking
%   F09_nlp_bert.png                     →  fig:nlp_bert
%   F10_master_summary.png               →  fig:master_summary
%
% All files are at:
%   /Volumes/Research/Subho_IIM/Robust-NN-learning/results/paper/figures/
% ===========================================================================

\begin{figure}[p]
  \centering
  \includegraphics[width=\linewidth]{figures/F01_clean_accuracy_all_datasets.png}
  \caption{Clean-data test accuracy (\%) of all nine loss functions across
  MNIST, CIFAR-10, PathMNIST, and DermaMNIST ($\eta{=}0$, seed~42).
  SDIV is competitive on MNIST and PathMNIST; on DermaMNIST the default
  parameters trigger majority-class collapse (66.88\%, dashed line).}
  \label{fig:clean_all}
\end{figure}

\begin{figure}[p]
  \centering
  \includegraphics[width=\linewidth]{figures/F04_noise_mnist.png}
  \caption{Test accuracy on \textbf{MNIST} under uniform label noise
  $\eta \in \{0, 10, 20, 30, 40\}\%$ (seed~42).
  SDIV and GCE degrade the least ($-1.5$~pp and $-1.2$~pp respectively at
  $\eta{=}40\%$); CCE degrades $-3.6$~pp.}
  \label{fig:noise_mnist}
\end{figure}

\begin{figure}[p]
  \centering
  \includegraphics[width=\linewidth]{figures/F02_noise_pathmnist.png}
  \caption{Test accuracy on \textbf{PathMNIST} under uniform label noise.
  Near-flat trajectories for most losses indicate a ceiling effect.
  MAE collapses to 48.76\% at $\eta{=}10\%$ before partially recovering
  (AUNRC = 0.2109, vs CCE = 0.3269).}
  \label{fig:noise_pathmnist}
\end{figure}

\begin{figure}[p]
  \centering
  \includegraphics[width=\linewidth]{figures/F03_noise_dermamnist_collapse.png}
  \caption{Test accuracy on \textbf{DermaMNIST} under uniform label noise.
  SDIV, MAE, and GCE show flat lines at $\approx$66.88\%~---
  the majority-class prevalence~--- not robustness but gradient collapse.
  FCL (70.72\% mean) and TSCCE (70.47\% mean) maintain the most stable
  trajectories among discriminative losses.}
  \label{fig:noise_dermamnist}
\end{figure}

\begin{figure}[p]
  \centering
  \includegraphics[width=\linewidth]{figures/F08_aunrc_ranking.png}
  \caption{Area Under the Noise-Robustness Curve (AUNRC) ranking on
  PathMNIST (trapezoidal integral of accuracy over $\eta \in [0,0.4]$).
  FCL (0.3321) $>$ SCE (0.3313) $>$ TPDD-CCE (0.3293) $>$
  SDIV (0.3279) $>$ CCE (0.3269).
  MAE (0.2109) is an extreme outlier due to early collapse.}
  \label{fig:aunrc_ranking}
\end{figure}

\begin{figure}[p]
  \centering
  \includegraphics[width=\linewidth]{figures/F05_fgsm_pathmnist.png}
  \caption{FGSM adversarial accuracy on \textbf{PathMNIST}
  ($\varepsilon \in \{0, 1/255, 2/255, 4/255, 8/255\}$, seed~42).
  All losses collapse catastrophically at $\varepsilon{=}8/255$;
  no robust training objective confers genuine adversarial resistance.}
  \label{fig:fgsm_pathmnist}
\end{figure}

\begin{figure}[p]
  \centering
  \includegraphics[width=\linewidth]{figures/F06_fgsm_dermamnist_artefact.png}
  \caption{FGSM adversarial accuracy on \textbf{DermaMNIST}.
  SDIV/MAE/GCE appear immune (flat at $\approx$66.88\%)~---
  a majority-class collapse artefact, not genuine adversarial robustness.
  CCE falls 49.6~pp; SCE ($-17.1$~pp) and TSCCE ($-20.6$~pp) degrade
  substantially less among discriminative losses.}
  \label{fig:fgsm_dermamnist}
\end{figure}

\begin{figure}[p]
  \centering
  \includegraphics[width=\linewidth]{figures/F07_sdiv_grid.png}
  \caption{SDIV $(\beta, \lambda)$ parameter response surface on
  PathMNIST (left) and DermaMNIST (right), evaluated at 17 valid
  grid points (seed~42, 30 epochs).
  Moderately negative $\lambda{=}{-0.40}$ consistently outperforms
  the default $\lambda{=}{-0.80}$;
  best PathMNIST: $(\beta,\lambda){=}(0.05,-0.40)$ at 84.11\%;
  best DermaMNIST: $(0.10,-0.40)$ at 73.32\%.}
  \label{fig:sdiv_grid}
\end{figure}

\begin{figure}[p]
  \centering
  \includegraphics[width=\linewidth]{figures/F09_nlp_bert.png}
  \caption{BERT fine-tuning accuracy on Emotion (6-class) and PubMedQA
  (3-class) across nine loss functions (seed~42, 1{,}500 training samples,
  3 epochs). Differences $\leq 2.7$~pp; results are indicative only.}
  \label{fig:nlp_bert}
\end{figure}

\begin{figure}[p]
  \centering
  \includegraphics[width=\linewidth]{figures/F10_master_summary.png}
  \caption{Master summary: best test accuracy per loss function and per
  dataset across all four evaluation categories.
  SDIV (highlighted) is competitive with the field on standard benchmarks
  and, after parameter tuning, matches CCE on both medical imaging datasets.}
  \label{fig:master_summary}
\end{figure}
